The methodological challenges documented in \autoref{sec:discussion} translate into six concrete directions for future work: standardised reporting and open data, causal mechanisms, generalisation and ecological validation, sex- and gender-stratified design, objective behavioural markers, and longitudinal characterisation of adaptation and recovery.

First, standardised reporting represents a gap documented across the corpus.
For classification and prediction studies, the minimum items follow from the gaps documented in \autoref{sec:discussion}: a discriminative metric (AUC or balanced accuracy) with 95\% confidence interval, sensitivity and specificity at a stated operating point, the numeric label threshold, a subject-respecting validation scheme, and comparison against a trivial and a non-neural feature baseline.
For studies that quantify in-exposure neural change, a baseline matched to the exposure condition on visual input, oculomotor state, and vigilance is required for reproducibility, with the choice and its rationale reported; the available options are discussed in \autoref{sec:discussion}.
For EEG, the reference electrode and notch-filter frequency are required for reproducibility and meta-analytic re-pooling.
For subjective sickness measurement, reporting the questionnaire instrument, its administration time-point relative to exposure end, and the numeric cutoff used to assign class labels enables cross-study comparability; where continuous in-task ratings are used, the response modality and any associated motor artefact exclusion strategy are necessary for reproducibility.
For stimulus and display, the device model, field-of-view, and frame rate alongside the stimulus class (passive or active, recorded or interactive); adoption of a small number of canonical reference paradigms, at minimum a passive visual-flow exposure \cite{suwandiExploringImpactMotion2024} and an active first-person navigation task \cite{cortesEEGbasedExperimentVR2023a}, would allow direct cross-laboratory comparison without the confound of idiosyncratic stimulus design.

These reporting standards are a precondition for, but not a substitute for, shared infrastructure: community-curated benchmark datasets with multimodal recordings (EEG, EOG, ECG, head motion, sway, eye tracking), well-defined subjective labels, and coverage of both passive and active paradigms.
A small number of corpus papers have already released public datasets \cite{demirelSubjectivityContinuousCybersickness2025, kimDeepCybersicknessPredictor2019, dasdemirClassificationEmotionalImmersive2023}; releasing analysis code and pre-trained models alongside the data would further support replication and cumulative analysis.

Second, the causal status of documented neural associations remains largely unexamined.
The literature has documented associations between specific neural patterns and VIMS across multiple studies, but whether these associations are causal has not been established.
Non-invasive brain stimulation offers a direct test: transcranial alternating-current stimulation can target specific oscillations (e.g., alpha or theta), and existing corpus studies that combine stimulation with VR exposure \cite{takeuchiModulationExcitabilityTemporoparietal2018, bergerInfluencePlaceboTDCS2024} provide a methodological template for asking whether modulating a candidate neuromarker changes sickness severity.
Coupled with computational models of multisensory integration and predictive coding \cite{nurnbergerMismatchVisualVestibularInformation2021}, interventional designs can begin to distinguish neural features that drive the VIMS experience from features that merely accompany it.

Third, generalisation and ecological validation remain open challenges for models built for VIMS.
In the 4 corpus studies that report both within- and cross-subject accuracy, cross-subject performance is lower in each case, and only 1 of 38 classification studies reports a leave-one-subject-out evaluation; transfer-learning and domain-adaptation approaches that explicitly model inter-subject and inter-session variability address a current gap if brain-directed methods are to operate without per-user calibration.
Validation beyond the laboratory is a further open direction: at-home, naturalistic paradigms such as the in-home gaming setup of \cite{moinnereauInstrumentingVirtualReality2022} address the question of whether laboratory neuromarkers survive the noise of consumer-grade hardware and uncontrolled environments.

Fourth, sex and gender represent an underexplored factor in the corpus.
Of the 92 corpus studies, 11 used male-only samples, 15 did not report sex composition, and only 3 (all from a single laboratory) reported higher cybersickness in women \cite{bergerSexDifferencesUser2021, bergerInfluencePlaceboTDCS2024, bergerManipulatingCybersicknessVirtual2025}.
The broader literature on sex differences in VIMS susceptibility is mixed: \cite{kellyGenderDifferencesCybersickness2023} argues that reported effects are modest in magnitude and plausibly mediated by interpupillary-distance fit, prior gaming experience, and questionnaire framing; \cite{kourtesisCybersicknessCognitionMotor2023} reports that male--female differences disappear when prior gaming experience is included as a covariate; and \cite{lawsonAttentionMilitaryCommercial2023} caution that the claim of greater female susceptibility is overstated in the simulation literature.
Taken together, the 3 corpus reports and the broader literature are consistent with a modest, variable effect whose magnitude and moderators cannot be resolved at individual-study scale, and whose apparent direction may partly reflect gaming-experience differences and questionnaire framing rather than a stable sex difference in susceptibility.
Reporting sex-stratified outcomes (band-power changes, connectivity, classification metrics), even when single-study samples are underpowered for inferential testing, would enable subsequent meta-analytic re-pooling.
Logging interpupillary-distance setting and gaming experience as covariates would address both the sex confound and the hardware effects documented in \autoref{sec:confounders}.
Adequately powered designs, such as pooled multi-site cohorts or longitudinal within-subject designs, are the precondition for resolving the magnitude and moderators of a sex effect on VIMS.

Fifth, objective behavioural markers represent an underdeveloped line of investigation.
The candidate markers introduced in \autoref{sec:discussion} are already collected by 16 of the 92 corpus studies, but typically as covariates rather than as candidate ground truth.
Postural sway is the strongest theoretically anchored candidate: reporting a coupling estimate between pre-symptomatic sway dynamics and subsequent self-reported sickness, alongside the protocol metadata that this estimate depends on (stance, surface, sampling rate), would anchor the construct across laboratories.
Pupillometry can supplement self-report subject to arousal and cognitive-load confounds \cite{kourtesisCybersicknessVirtualReality2023}, and requires controls on ambient luminance and on task demand.
Heart-rate variability, electrodermal activity, and head-motion dynamics are further candidates whose corpus-level predictive value is not yet quantified \cite{islamAutomaticDetectionPrediction2020, tianWhoSaysYou2024, yangPredictionDetectionVirtual2023}; multi-modal integration that fuses behavioural, autonomic, and neural signals offers a more holistic state estimate than any single modality.
Treating these markers as outcome variables, rather than as covariates, is the precondition for shifting VIMS research toward objective ground truth.

Finally, longitudinal designs represent a gap in the current corpus.
The corpus is dominated by single-session experiments, leaving the neurodynamics by which users adapt to VR (``getting one's VR legs'') largely uncharacterised, and the single corpus recovery study \cite{wooRecoveryTimeVR2023} is insufficient to support strong claims about residual neural state.
Multi-session within-subject designs would enable tracking of how the EEG signatures identified in this review evolve with repeated exposure, and mapping the dissipation of post-exposure signatures onto behaviour.
